In this thesis, we presented three studies pertaining to diverse, yet related facets of pandemic preparedness: The cornerstone of this thesis revolves around engineering new data formats and analysing of lesser explored datasets which may complement existing approaches to understanding the evolution and transmission of pathogenic diseases with epidemic potential. 

On the one hand, the overarching goal of phylo2vec was to provide a new bridge between the worlds of phylogenetics and machine learning through a compact vectorization of phylogenetic trees. Drawing inspiration from tree enumeration strategies~\cite{rohlf1983}, this study has breathed new life into the development of new vector-based representations of phylogenies such as HOP~\cite{chauve2025} and OLA~\cite{richman2025}.

% Currently limited scope beyond its compression/efficient representation/fast sampling, which are cool but not detrimental problems in phylogenetics. Could be a no free lunch thing, where operations are more difficult to imagine when you start from a vector that describes a branching process vs. an explicit tree itself. Although the new edition of the software in Rust/Python/R has more functionalities, the next step is likely to really dabble further into phylogenetic inference. Another possibility is data visualisation, but that would be a bit far from the original advantages of phylo2vec.

Nevertheless, the path to their usage in machine or deep learning applications relevant to epidemiology remains uncertain. In particular, it is unclear how an end-to-end deep neural network-based solution for general phylogenetic inference could be constructed~\cite{sapoval2022}. At its essence, deep learning requires a neural network and data from which the network learns from, and a metric to evaluate the learning performance. As an example, consider maximum likelihood-based molecular phylogenetic inference as a task to model with deep learning. Input data typically include an alignment of $n$ molecular sequence with $m$ sites and a model of evolution. Outputs are a single bifurcating tree and associated parameters which maximize Felsenstein's likelihood~\cite{felsenstein1973, felsenstein1981}. Although this metric may be adopted for deep learning-based applications, it is not fully differentiable due to the discreteness of tree topology, at least not without continuous relaxations~\cite{penn2023} or projections in non-Euclidean spaces~\cite{mimori2023, macaulay2024}. Furthermore, current approaches involving gradient computations on trees are not yet as efficient as precisely handcrafted approaches~\cite{fourment2023}. In addition, these approaches have not yet managed to fully leverage \glspl{gpu} to parallelize computations~\cite{fourment2023}, which is a hallmark of the success of deep learning~\cite{krizhevsky2012, lecun2015, vaswani2017}. Beyond issues related to tree likelihood, building comprehensive and biologically plausible training datasets is also non-trivial. In general, the validity of an inferred tree is inherently constrained by available data and models of evolution~\cite{yang2014}. In more specific cases involving reticulate processes (e.g., recombination), binary trees may provide an incomplete picture of evolutionary history~\cite{huson2006, wong2024}. To that end, several deep learning-based approaches have made use of simulated data from trees, but are yet to outperform likelihood-bases approaches~\cite{nesterenko2025}. Last, but most prominently, a deep learning approach would need to be both sufficiently flexible and scalable to incorporate a hugely variable-length input and predict an optimal tree from a combinatorially vast tree space~\cite{sapoval2022}. Various sequence representation approaches have been explored, including chaos game representations~\cite{jeffrey1990} and learnable representations from Transformer-based models~\cite{consens2025}. However, these approaches are yet to provide additional insight on evolution~\cite{reasoning}. Alternatively, graph-based generative models may also have potential for generating plausible trees~\cite{xie2023, xie2025}, but currently remain limited for practical phylogenetic applications. Consequently, an emerging body of research at the intersection of deep learning and phylogenetics has focused on optimising model parameters through tools like ModelTeller~\cite{abadi2020}, ModelFinder~\cite{kalyaanamoorthy2017}, and PhyloDeep~\cite{voznica2022}, which constitutes a promising avenue for phylo2vec usage.

Finally, the development of phylo2vec also motivated the development of a high-performance software for phylogenetic tree manipulation~\cite{scheidwasser2025}, comprising various functions for tree sampling, format conversions and operations on trees. However, key features are still lacking for the software to become a fully featured and widely adopted tool in the bioinformatics community. While the computational efficiency of the library may facilitate rapid tree visualisation, the non-equivalence of topology moves in the phylo2vec space to rooted \gls{spr} moves represents a promising avenue for the development of phylogenetic inference based on phylo2vec.

% Key features are still missing to make the software fully-fledged and attractive to the bioinformatics community, notably phylogenetic inference. 

% the super-exponential nature of tree space makes scalability a considerable challenge for deep learning-based solutions.

%  % First, ``true" evolutionary trees may be difficult or impossible to establish. For some loci or genes involving reticulate processes (e.g., recombination), binary trees may provide an incomplete picture of evolutionary history~\cite{huson2006, wong2024}. 

% Whereas 

% may be difficult or impossible to definitively establish, particularly when  are involved. simualtions

% % it is not immediately differentiable due to its recursive nature.

% However, building a comprehensive and plausible training dataset is nontrivial. \textcolor{red}{True trees may be difficult if not impossible to infer, e.g. in the case of reticulate processes (tree mixtures and phylo networks might be better). Second, what is a true tree anyway? Third, model of evolution?}. In addition, a deep learning approach would need to be both sufficiently flexible and scalable to incorporate a hugely variable-length input and predict a single tree from a potentially exponentially vast tree space. \textcolor{red}{chaos game representation for sequences~\cite{jeffrey1990}, representing sequences as text (hydra)}. \textcolor{red}{graphvae, graphrnn, diffusion have potential to generate trees/graphs (phylovae?)}, but they are poorly scalable. \textcolor{red}{phyloformer does not beat iqtree, and is worse as the number of sequences grows. gradme = not deep learning but gradient based, not better than iqtree.}. Thus, a body of research at the intersection of deep learning and phylogenetics is consacred to optimizing model parameters: ModelTeller, PhyloDeep~\cite{voznica2022} etc. This also constitutes a promising avenue for phylo2vec usage. \textcolor{red}{I have thought of MCTS for phylo2vec inference, but scalability remains an issue}.

The second study of this thesis aimed at exploring the potential added value of data from protein structures into giving insights on viral/pathogenic evolution. Mixed bag of results, difficult problem. Foldseek trained on a cluster of proteins and may not be fine-grained enough to accurately model all angles? 
% In molecular phylogenetic inference, input data typically consists of $N$ molecular sequences and the output typically consists of an optimal bifurcating tree (and associated parameters), or a distribution of trees in Bayesian phylogenetics.

% \textcolor{red}{chaos game representation}  with deterministic ...  On the one hand, 

% collecting a valid and comprehensive training dataset is a nontrivial task, as obtaining a ``true" bifurcating phylogeny may not be possible. For instance, reticulate processes of evolution such as horizontal gene transfer or recombination cannot accurately be modelled by inferring a single bifurcating tree. Instead, approaches based on mixtures of trees and phylogenetic networks have proved to be more insightful and plausible~\cite{huson2006, wong2024}. 

% is practically impossible, as obtaining a ``true" bifurcating phylogeny may not be possible

% A general deep learning solution would have to face face  due to scalability (tree space grows exponentially) and lack of training data (obtaining a ``true" bifurcating phylogeny is   promising avenue is the inference of model parameters

% Inspired by other tree labelling strategies~\cite{rohlf1983, dinh2010}

% exploring how lesser explored representations of pathogenic data may h

Summary of the results of the papers and their relation to international state-of-the-art research within the subject area
If required for the studies, information on ethical and legal permits and approvals
Conclusions and perspectives for further research

Despite claims that "we could AlphaFold our way through the next pandemic", I'm not sure. Folding tons of proteins remains expensive, whether with deep learning or x-ray crystallography. Like a lot of things with AI, it's easy to get 90\% there, but the real work is the last 10\%. For pandemic data, this is also the case; currently it's still a doubt whether these systems can help with variant-level/individual-level structures. Whereas there is undoubtedly some value in understanding the structure of antigens, its increased complexity compared to sequences might not provide significant benefits compared to sequence-based phylogenetics. Current research should focus on exploiting "rapid response to fast viral evolution with topology-assisted..." or fine-tuning either protein structure models or structure-discretizers such as 3di to single viruses. All in all, this is again a sign that there is not single tool to rule all pandemic problems.

From a scientific point of view, pandemic preparedness is hugely multi-faceted. And within each of these facets, the skill sets needed are also diverse. Take chicken study for instance: we have deep learning, videography, microbiology, veterinary science.

Video-based animal surveillance, although a nice concept, is incredibly difficult due to the number of degrees of freedom.

While optical flow is limited in resolution/precision, deep learning-based tracking is not yet ready for farm-scale/bird-scale environments, where 100/1000+ animals coexist. It has lots of potential, both from a public health and an economic perspective, but scability is currently the Achille's heel. My personal idea would be a two-camera system: one cheap one, where we could perform basic optical flow surveillance. If there are zones with (sustained) lower-than-average inactivity, a second camera hooked to raspberry pi + GPU for tracking tries to assess the welfare of the individual in this zone. This would be a "best of both worlds" strategy: we use cheap optical-flow as it's fast, and restrict the number of animals to track using deep learning. A robot is also envisageble, but that poses additional logistical issues + invasive.

This PhD illustrates the need for collaboration between all sorts of sciences, from agricultural scientists and veterinarians to lab technicians to computer scientists. For vets, there is a need to develop and openly share such data to advance the field. Problematic in industries due to ethical concerns + whistleblowing?

Vectorising pandemic data is in any case important because of the sheer volume of data we collect. Terabytes of video data for animal surveillance, terabytes of sequences for large-scale human genomic surveillance, terabytes of protein structures. Being able to efficiently compress data is green (less memory, less heat? etc.), and makes sharing data easier.